Embora intimamente relacionado às funções, o escopo de variáveis é um 
conceito que merece atenção própria. Em termos gerais, o escopo refere-se 
a um conjunto de regras que determinam como as variáveis podem ser usadas em 
um programa.

\subsubsection{Variáveis locais}

Variáveis locais são aquelas que só podem ser acessadas em uma função ou bloco 
de código específico. Em outras palavras, uma vez declaradas numa função, elas 
não podem ser utilizadas ou modificadas em outros contextos, e deixam de 
existir ao término da função.

Para ilustrar essas propriedades das variáveis locais, considere o seguinte 
exemplo:

\begin{code}
    #include <stdio.h>

    void printnum(int x) { // Recebe argumento 'x'
    printf("\%i\n", x);

    int y = 369; // Cria variável local
    printf("\%i\n", y);
    }

    int main (void) {
        printnum(369);
        return 0;
    }
\end{code}

Neste código, a função \texttt{printnum} imprime dois valores na tela: o 
argumento \(x\) passado para a função e a variável local \(y\) declarada 
dentro dela. Ambos os valores são considerados variáveis locais, pois só 
podem ser acessados e modificados dentro da função em que foram declarados.

\begin{ps}
    Em alguns contextos, as variáveis locais recebidas como argumentos são 
    chamadas de ``parâmetros locais'', mas, para fins de entendimento e 
    praticidade, não há necessidade de distinção, uma vez que ambas são 
    manipuladas da mesma forma.
\end{ps}

\subsubsection{Variáveis globais}

Similarmente às variáveis locais, as variáveis globais podem ser acessadas e 
modificadas por funções ou blocos de código. No entanto, elas têm a 
característica de serem acessíveis de qualquer parte do código e permanecerem 
em existência após o término de um bloco ou função.

Em aplicações simples, variáveis globais podem ser extremamente úteis, pois 
eliminam a necessidade de passar valores entre funções para realizar 
determinadas tarefas. No entanto, em aplicações mais complexas, o uso 
excessivo de variáveis globais pode levar a problemas de controle de acesso, 
já que diferentes partes do programa podem tentar acessar ou modificar a mesma 
variável simultaneamente, o que pode acabar gerando algumas inconveniências.

Apesar de inicialmente intimidador, o uso de variáveis globais é uma prática 
que deve ser dominada, pois é fundamental para compreender o fluxo de dados 
em programas mais complexos.

Considere o seguinte exemplo prático de como uma variável global pode ser 
utilizada em um projeto real, como para controlar a vida do personagem em um 
jogo:

\begin{code}
    #include <stdio.h>

    unsigned int life = 100; // Vida do personagem

    // Danifica a vida do personagem
    void lifeDamage(int points) {
        life -= points;
    }

    // Recupera a vida do personagem
    void lifeHeal(int points) {
        life += points;
    }

    // Exibe a vida do personagem
    void lifePrint(void) {
        printf("Vida: \%i\n", life);
    }

    int main(void) {
        lifePrint();

        lifeDamage(10);
        lifePrint();

        lifeHeal(10);
        lifePrint();

        // Modifica o valor da vida diretamente
        life = 10;
        lifePrint();

        return 0;
    }
\end{code}

Neste código, a variável \texttt{life} é uma variável global que armazena 
a vida do personagem. Ela pode ser acessada e modificada por todas as 
funções do programa, que simplifica o controle do estado do personagem ao 
longo do jogo. É notável também que, ao contrário das variáveis locais, uma 
variável global mantém seu valor durante toda a execução do programa.

\subsubsection{A diretiva static}

Variáveis precedidas pela diretiva \texttt{static} mantêm seus valores mesmo 
quando saem de escopo. Isso significa que, ao executar uma função que contém 
uma variável \texttt{static}, a variável é inicializada apenas na primeira vez 
que a função é chamada, e seu valor é preservado para chamadas subsequentes.

Variáveis \texttt{static} são alocadas em uma área da memória do programa 
conhecida como \texttt{data}, enquanto todas as outras variáveis que 
aprendemos a declarar até agora são alocadas em uma área conhecida como 
\texttt{stack}. Por isso, as variáveis \texttt{static} são inicializadas com o 
valor \(0\) por padrão, a menos que o usuário as inicialize explicitamente com 
um valor diferente.

O código a seguir ilustra essa propriedade:

\begin{code}
    #include <stdio.h>

    int main(void)
    {
        static int x;
        int y;
        printf("x: \%d\n", x);
        printf("y: \%d\n", y);
    }
\end{code}

Ao compilar e executar o código acima você vai notar que a variável \(y\) 
terá um valor aleatório, enquanto a variável \(x\) sempre terá o valor 0. 
Isso ocorre porque, quando uma variável é alocada no espaço \texttt{data} da 
memória, seu valor deve ser constante.

Usemos esse código como exemplo:

\begin{code}
    int initializer(void)
    {
        return 50;
    }

    int main(void)
    {
        // Inicializa o valor de x com o valor retornado pela função
        static int x = initializer(); 
        printf("x: \%d\n", x);

        return 0;
    }
\end{code}

Ao compilar o código acima você verá um erro semelhante a:

\begin{code}
    main.c: In function ‘main’:
    main.c:10:24: error: initializer element is not constant
    10 |         static int x = initializer();
    |                        ^~~~~~~~~~~
\end{code}

Isso se deve ao fato de que o valor de \(x\) não é conhecido durante a 
compilação do programa, sendo assim dependente de outra função ou variável 
para ser atribuído.

\begin{ps}
    Em relação aos espaços \texttt{heap} e \texttt{data} da memória, esses 
    conceitos serão discutidos com mais detalhes no capítulo \ref{secoes-da-memoria}
\end{ps}

Além disso, é possível declarar variáveis \texttt{static} no contexto de 
variáveis globais e funções. Quando isso é feito, o compilador limita o uso 
dessa variável (ou função) ao arquivo em que foi declarada. Por exemplo, em 
programas mais complexos que são divididos em vários arquivos com diferentes 
segmentos de código, pode haver variáveis globais e funções com o mesmo nome, 
mas com propósitos diferentes. Nesses casos, a diretiva \texttt{static} pode 
ser útil.

O código a seguir ilustra o uso de variáveis globais \texttt{static}:

\begin{code}
    // Arquivo: main.c
    #include <stdio.h>

    static int x = 10;

    int main(void)
    {
        printf("x in main.c: \%d\n", x);
        printX();
        return 0;
    }

    // Arquivo: example.c
    #include <stdio.h>

    static int x = 20;

    void printX(void)
    {
        printf("x in example.c: \%d\n", x);
    }
\end{code}

E aqui está um exemplo de código que ilustra o uso de funções \texttt{static}:

\begin{code}
    // Arquivo: main.c
    #include <stdio.h>

    static void printHello(void)
    {
        printf("Hello from main.c!\n");
    }

    int main(void)
    {
        printHello();
        return 0;
    }

    // Arquivo: example.c
    #include <stdio.h>

    static void printHello(void)
    {
        printf("Hello from example.c!\n");
    }
\end{code}

\subsubsection{Exercícios}

\textbf{É de extrema importância que você realize os exercícios aqui elaborados para fixar o conteúdo.}

\begin{enumerate}
    \item {Escreva um programa que declare duas variáveis com o mesmo 
        nome, uma global e uma local em uma função que não seja a 
        \texttt{main}. O que acontece quando você tenta imprimir o valor 
        dessas variáveis dentro da função? 
        E fora da função?}
\end{enumerate}
